\section{Arithmetic compatibility of the actual choices}
\label{sec:arithmetic-compatibility}

This section supplies the rank and character comparisons upstream of NC.
Throughout, $D$ is one fixed datum, $E=D.\mathrm{host}.E$,
$L=D.\mathrm{what}_2.LL$, $L_0=D.LL0$, and
$\iota_L:L\hookrightarrow\mathbb C$ is its specified embedding.
The notation $\bar v$ denotes the specified conjugation on coefficient
places, while conjugation of primes of $E$ is written $\bar P$.
The actual coefficient rings, period modules and realization maps of $D$
are retained. In particular, a new auxiliary row model cannot be substituted
for the given $D$ when applying a rank comparison.

Write $G=G_E$, $G_f=U(E)(\mathbb A_f)$, and $\Gamma=U(E)(E)$ in the
source's standard positive definite unitary realization. For a twist index
$i$, put $K_{i,k}=\operatorname{Frac}(O_{i,k})$. A reading $\rho$ has an
injective map $\rho_L:L\to A_\rho=\operatorname{Frac}(O_\rho)$ and an
injective map $d_\rho:O_{i,k}\to A_\rho$. The expression
$\rho.\mathrm{Pulls}(\psi,\psi_0)$ means the literal equality
\[
 \rho_L(\psi_0(g))=d_\rho(\psi(\operatorname{art}_{i,Q}(g)))
 \qquad(g\in G).                                      \tag{RC.1}
\]
The normalized packet is
$p(i,b,Q,k,\psi)=(i,Q,k,\psi^{\epsilon_b})$, with
$\epsilon_{\mathrm{true}}=1$ and
$\epsilon_{\mathrm{false}}=-1$. Applying this side operation twice gives
$\psi$ itself. All Frobenius conventions below are those of the source:
an arithmetic Frobenius $\sigma_P$ has
$\chi(P)=r(\sigma_P)^{-1}$.

\paragraph{RC1: finite evaluations prove scalar extension for the form space.}
Let $K\hookrightarrow\mathbb C$ be an embedding $e$. Let $W_K,W_C$ be
finite-dimensional representations of $\Gamma$, with bases on the same
finite index set, such that every matrix of $\rho_C(\gamma)$ is the
entrywise image under $e$ of the corresponding matrix of $\rho_K(\gamma)$.
This is the exact predicate \texttt{WtMatch}; equality of dimensions alone
is insufficient. Let $J\le G_f$ and $\chi:J\to K^\times$. Assume that
$J$ contains a positive congruence subgroup on which $\chi=1$.
For
\[
 \mathcal A_K=\{F:G_f\to W_K:
 F(\gamma g)=\rho_K(\gamma)F(g),\
 F(g\kappa)=\chi(\kappa)F(g)\},
\]
assume $\dim_K\mathcal A_K<\infty$. Then the same equations over $\mathbb C$,
with character $e\chi$, define precisely
$\mathbb C\otimes_{K,e}\mathcal A_K$. In particular,
\[
 \dim_K\mathcal A_K=\dim_{\mathbb C}\mathcal A_C.          \tag{RC.2}
\]

First, any finite-dimensional subspace $V$ of functions $S\to W$ is
detected by finitely many evaluations. Starting with $V$, if the current
common kernel is nonzero, choose a nonzero function in it and a point where
it is nonzero. Adding that evaluation strictly decreases the kernel
dimension. After at most $\dim V$ steps the kernel is zero.
Choose such a finite set $S_0$ for $\mathcal A_K$.

The common weight bases identify the coefficientwise map
$W_K\to W_C$ and intertwine every $\Gamma$-action. Images of a $K$-basis
of $\mathcal A_K$ are $\mathbb C$-independent: evaluation on $S_0$, followed
by the finite weight-coordinate map, gives a finite matrix of full column
rank over $K$; a nonzero maximal minor remains nonzero under $e$.

For spanning, choose a basis $(\beta_a)$ of $\mathbb C$ over $e(K)$.
For $F\in\mathcal A_C$, expand each weight coordinate in this basis.
Extracting the coefficient of $\beta_a$ gives a $K$-valued function $F_a$.
Because the representation matrices and the level character are defined
over $e(K)$, each $F_a$ satisfies both defining equations and lies in
$\mathcal A_K$. The finitely many coordinates of $F|_{S_0}$ involve only
finitely many $\beta_a$, say those in $A_0$. If $a\notin A_0$, $F_a$ vanishes
on $S_0$, hence everywhere. Thus
$F=\sum_{a\in A_0}\beta_a e(F_a)$ globally, and each $F_a$ is a finite
$K$-combination of the chosen form basis. This proves spanning.
The congruence-kernel hypothesis supplies the smoothness required in
\texttt{NG.TypeSp}: every function with the stated right transformation
law is fixed by that same congruence subgroup. No commutation of scalar
extension with an arbitrary infinite product has been assumed.

\paragraph{RC2: the literal level and weight comparison supplies RankTie.}
Let $m:\texttt{MtReal}(D,\mu,j')$ and $X:\texttt{NExpTied}(D)$.
For an admissible twist $(i,b,k,Q,\psi)$ and $j\in\operatorname{Fin}(D.R.m)$,
put $c=X.\mathrm{mtV}(i,j+1,b,k,Q,\psi)$.
Suppose $e:K_{i,k}\hookrightarrow\mathbb C$ satisfies
\[
 \texttt{LevEq}(m.J,m.\chi,c.J,e\,c.\chi),\qquad
 \texttt{WtMatch}(e,c.\rho,m.\mathrm{wt}.\rho).
                                                               \tag{RC.3}
\]
Then
\[
 \dim_L \mathrm{MtL}(\mu,j')
 =\dim_{K_{i,k}}\mathrm{MtV}(i,j+1,b,k,Q,\psi).             \tag{RC.4}
\]
The local lattice is finitely generated and spans the actual local
fraction-field module, so that module is finite-dimensional.
The map $c.\mathrm{real}$ identifies it with its $K_{i,k}$-form space.
RC1 changes coefficients to $\mathbb C$. The first predicate in (RC.3)
is equality of subgroups and of their character values, so it identifies
the resulting complex form space with the range of $m.\mathrm{real}$.
The latter map is injective and its range is exactly that form space;
therefore its dimension is $\dim_{\mathbb C}\mathrm{Mt}(\mu,j')$.
Finally, the actual \texttt{jMt\_basis} field says that the images of an
$L$-basis under $\mathrm{jMt}$ form a complex basis. This gives (RC.4).

Apply this to $\mu=\mu_{\mathrm{Tw}}(v)$ and $j'=j+1$ for every
$1\le j'\le D.R.m$. It gives \texttt{RankEq} for the unnormalized
character and \texttt{RankTie} for its normalized packet, because the
two applications of $\psi\mapsto\psi^{\epsilon_b}$ cancel.
Neither an abstract dimension estimate nor the existence of an unrelated
row model supplies (RC.3).

\paragraph{RC3: finite-order Hecke characters have the required coefficient-valued Artin character.}
Retain the exact published global reciprocity input: a surjective open
quotient homomorphism
\[
 \operatorname{Art}:\mathbb A_E^\times\longrightarrow G_E^{\mathrm{ab}},
\]
whose kernel is the closure of principal ideles times the connected
infinite component, and whose inertia, uniformizer and decomposition
images are the ones stated in \texttt{IsGlobalArtin}.
Let $\chi$ be a continuous Hecke character with $\chi^n=1$, $n>0$,
and suppose $L$ contains a primitive $n$th root of unity.
There is a character $\psi_0:G_E\to L^\times$, with open kernel and
$\psi_0^n=1$, whose complex image corresponds to $\chi$ under this
specific Artin map. Thus the literal \texttt{GalOf} predicate holds.

The image of $\chi$ is contained in the finite set $\mu_n(\mathbb C)$.
Its kernel is both open and closed. Principal ideles are killed by the
Hecke-character law, and the connected infinite component has connected
image in a finite set, hence is also killed. Closedness of the kernel
therefore kills the entire Artin kernel. Quotient descent gives a
homomorphism $G_E^{\mathrm{ab}}\to\mu_n(\mathbb C)$; the open quotient
property proves its continuity. Compose with abelianization.
The map $\mu_n(L)\to\mu_n(\mathbb C)$ induced by $\iota_L$ is bijective:
both groups consist of powers of the specified primitive root and its
image. Pulling back the character through this bijection gives $\psi_0$,
preserves its kernel, and proves all the assertions.

This construction can be chosen as a function of the pair $(\chi,n)$
alone. That choice is used below to obtain exact equalities of common
characters at places with the same data. It does not require independent
choices of a Galois character at every coefficient place.

\paragraph{RC4: Frobenius detection follows from the stated reciprocity map and weak approximation.}
Let $\varphi:G_E\to B$ be a continuous homomorphism to a Hausdorff
commutative topological group. Suppose that, outside a finite set $S$,
$\varphi(\sigma)=1$ for every valuation above $P$ and every arithmetic
Frobenius $\sigma$ at it. Then $\varphi=1$.

Descend $\varphi$ through topological abelianization and compose with
$\operatorname{Art}$, obtaining a continuous idele character $\theta$.
It is trivial on principal ideles. The uniformizer-image clause of
\texttt{IsGlobalArtin}, with its inverse-Frobenius convention, says that
$\theta$ kills every local uniformizer at each $P\notin S$.
If $\pi$ is one such uniformizer and $x\in E_P^\times$, then
$x\pi^{1-v_P(x)}$ is another; hence $\theta(x)=1$.
Thus every local multiplicative group outside $S$ is killed.

The resulting character depends only on the factors at $S$ and the
infinite places. To justify this assertion without an infinite product
calculation, approximate an idele that is $1$ at those factors by its
finite-coordinate truncations; these converge in the restricted-product
topology, and all are killed. Weak approximation makes the diagonal
$E^\times$ dense in the finite product
$\prod_{P\in S}E_P^\times\times E_\infty^\times$.
Indeed, additive weak approximation applies simultaneously to a finite
set of nonzero coordinates, and the unit locus is open, with continuous
inversion. The induced character on this product is therefore $1$.
Surjectivity of $\operatorname{Art}$ gives $\varphi=1$.
This argument uses the precise global reciprocity input of RC3; it does
not add a Chebotarev hypothesis.

\paragraph{RC5: the common character has exactly the required local behavior.}
Assume $\chi$ is an \texttt{IsPsi} character for $\ell$ and $Q$, and
$\psi_0$ satisfies \texttt{GalOf} as in RC3. Then $\psi_0$ is
anticyclotomic, is ramified exactly at $Q$, and is trivial on the entire
decomposition group above each prime over $\ell$.

Here the source host supplies the quadratic extension $E/F$, its
involution, and the unramified normalization
$v_P(a)=v_{P\cap F}(a)$ for $a\in F^\times$.
The condition $\chi|_{\mathbb A_F^\times}=1$ gives, at every pair of
unramified conjugate primes,
\[
                  \chi(P)\chi(\bar P)=1.                \tag{RC.5}
\]
At a split prime this is the product of the two local values on a
uniformizer of $F$; at a fixed prime the single local value is itself
$1$. The stated valuation normalization is what permits this uniformizer
calculation. Conjugation carries a valuation above $P$ to one above
$\bar P$, preserves the residue cardinality, and sends an arithmetic
Frobenius to an arithmetic Frobenius there. Therefore the continuous
character
$g\mapsto \iota_L(\psi_0(g)\psi_0(cgc^{-1}))$
is $1$ on every Frobenius outside a finite conjugation-stable set.
Its continuity follows from the open kernel of $\psi_0$ and continuity
of conjugation on the profinite Galois group; no continuity of a
coefficient-field isomorphism is needed. RC4 and injectivity of
$\iota_L$ give $\psi_0(cgc^{-1})=\psi_0(g)^{-1}$.

The inertia-image clause of reciprocity says that $\chi$ is unramified
at $P$ if and only if $\psi_0$ kills inertia there. The exact
\texttt{IsPsi.ram} biconditional therefore gives
\texttt{RamExactly}$(\psi_0,Q)$, in both directions.
At $P\mid\ell$, \texttt{IsPsi.atl} kills the whole local
$E_P^\times$. The closure of its Artin image is the abelianized
decomposition image. Since the kernel of the complex character is
closed, it kills that closure; injectivity then gives triviality of
$\psi_0$ on the decomposition group. Merely knowing unramifiedness at
$\ell$ would not prove this last statement.

\paragraph{RC6: actual ring-class descent gives all seven TwOK clauses.}
Fix a reading $\rho$ of $(i,b,k)$ at $v$ and a character $\psi_0$ with
open kernel, $\ell_i^k$-power exponent, anticyclotomic symmetry, exact
ramification $Q$, and decomposition triviality above $\ell_i$.
Assume that $Q$ consists of split places and is disjoint from $S_0(i)$.
Use the specified $\ell_i^k$th root $z_{i,k}\in O_{i,k}^\times$.
For every $g$, $\rho_L(\psi_0(g))$ is a power of $d_\rho(z_{i,k})$.
There is consequently a unique value
$\widetilde\psi(g)\in\langle z_{i,k}\rangle$ with
$d_\rho(\widetilde\psi(g))=\rho_L(\psi_0(g))$.
Injectivity of $d_\rho$ proves multiplicativity and equality of kernels.

Retain the actual \texttt{NTw.art\_univ} field. It factors every
open-kernel, anticyclotomic, finite $\ell_i$-primary character unramified
outside $Q$ through the specified finite group
$\operatorname{RCG}(i,Q)$, when $Q$ avoids primes over $\ell_i$.
That avoidance follows from $Q\cap S_0(i)=\varnothing$, since $S_0(i)$
contains all those primes. Applying the universal property to the finite
cyclic group $\langle z_{i,k}\rangle$ gives
\[
 \psi:\operatorname{RCG}(i,Q)\to O_{i,k}^{\times},
 \qquad \widetilde\psi=\psi\circ\operatorname{art}_{i,Q}.
                                                               \tag{RC.6}
\]
This is the actual descent, not an assumption of (RC.6).

Every descent satisfying (RC.1), not only the one just chosen, has the
seven \texttt{TwOK} properties: split support, disjointness from $S_0$,
open kernel, anticyclotomic symmetry, image in
$\langle z_{i,k}\rangle$, exact ramification $Q$, and triviality on
decomposition above $\ell_i$. The last five follow by applying the two
injective coefficient maps to the corresponding equalities for $\psi_0$.
The root condition follows because all roots of $X^{\ell_i^k}-1$ in
the fraction field are powers of $z_{i,k}$. These clauses imply the
literal \texttt{AdmTw}: $S_0$ also contains the type and auxiliary
exceptional primes, while the exponent and unramifiedness clauses follow
from the displayed cyclic image and exact ramification.
Existence of the ring-class object and its universal map remains part
of the retained \texttt{NTw} input.

\paragraph{RC7: two readings of the common character differ by an allowable power.}
Suppose two readings, possibly on different sides and at different
coefficient places, have the same $i,Q,k$, and their descents $\psi,\psi'$
read the same $\psi_0$. Suppose $L$ contains a primitive
$n=\ell_i^k$th root $z$ and $\psi_0^n=1$. Then
\[
 \psi'=\psi^a\quad\text{for some }a\in\mathbb Z
 \text{ with }\gcd(|a|,\ell_i)=1.                        \tag{RC.7}
\]
For $k=0$, both characters are $1$ by surjectivity of
$\operatorname{art}_{i,Q}$, and $a=1$ works.
For $k>0$, each image of $z$ and each image of $z_{i,k}$ is a primitive
$n$th root in the relevant reading fraction field. Thus changing from
one chosen primitive generator to the other is raising to an exponent
invertible modulo $n$. Write $\psi_0(g)=z^{t_g}$.
Comparing these powers in the two fraction fields gives
$\psi'(\operatorname{art}(g))
 =\psi(\operatorname{art}(g))^a$, with $a$ a product of two such
invertible exponents. Surjectivity of $\operatorname{art}$ gives (RC.7)
on the whole ring-class group.

For normalized packets the exponent is
$\epsilon_{b'}\epsilon_b a$; changing one side negates $a$.
It is still prime to $\ell_i$. This sign is essential: power ambiguity
is a statement about the normalized packet as well as about its
unnormalized character.

\paragraph{RC8: a reading supplies the exact side embedding and a complex extension.}
A reading determines
\[
 \rho_L(x)=\operatorname{Frac}(O_0(i)\to O_\rho)
                  (\operatorname{emb}_{i,b}(x))
       \quad(x\in L_0),                                 \tag{RC.8}
\]
and identifies the nonnegative-valuation predicate at $v$ with
integrality of $\operatorname{emb}_{i,b}(x)$ in $O_0(i)$.
To prove (RC.8), write the fractional image of $x$ with nonzero
denominator in $O_0(i)$ and apply the reading's literal
denominator identity. Fraction-field injectivity clears that denominator.
For the valuation assertion use the following reflection principle:
an injective local map from a valuation ring $A$ to a domain $B$
reflects integrality under its map of fraction fields.
If $t\notin A$, then $t^{-1}$ lies in the maximal ideal of $A$.
If the image of $t$ lay in $B$, the image of $t^{-1}$ would be a unit,
contrary to locality. The local maps needed here are supplied by NC4.
This proves the actual \texttt{SidePl} and \texttt{emb\_side} interfaces.

For a supplied \texttt{TwFrameC}, there is also an embedding
$\iota_\rho:A_\rho\hookrightarrow\mathbb C$ satisfying
$\iota_\rho\rho_L=\iota_L$ on all of $L$.
Here is the cardinality argument. A first-countable Hausdorff completion
of a dense set $S$ has cardinality at most $|S|^{\aleph_0}$: select one
convergent sequence for each point, and uniqueness of limits makes that
selection injective. Hence $|\mathbb Q_\ell|\le\mathfrak c$,
its algebraic closure has cardinality at most $\mathfrak c$, and its
completion $\mathbb C_\ell$ has cardinality at most
$\mathfrak c^{\aleph_0}=\mathfrak c$. The frame embeds $O_{i,k}$ into
$\mathbb C_\ell$. Finite generation of $O_\rho$ over $O_{i,k}$ and
formation of fractions therefore give $|A_\rho|\le\mathfrak c$.

Adjoin $\mathfrak c$ independent variables to $A_\rho$ and take an
algebraic closure. The resulting algebraically closed field has
characteristic zero and transcendence degree $\mathfrak c$ over the
number field $\rho_L(L)$. The same is true of $\mathbb C$ over
$\iota_L(L)$: an algebraic extension of a countable field on $\kappa$
indeterminates has cardinality $\max(\aleph_0,\kappa)$, so its
uncountable cardinality forces the stated transcendence degree.
Choose a bijection of transcendence bases extending
$\rho_L(x)\mapsto\iota_L(x)$, then extend the rational-function-field
isomorphism to algebraic closures and restrict to $A_\rho$.
This is an algebraic embedding; no analytic continuity is asserted.

\paragraph{RC9: coherent readings also give exact complex Hecke readings.}
Suppose families $s(v),Q(v),k(v),\psi_0(v)$ are unchanged under
$v\mapsto\bar v$, each $s(v)$ reads at $v$ on at least one side, and
every reading admits a descent of $\psi_0(v)$.
Assume $L_0/\mathbb Q$ is Galois, as required for the conjugate-reading
construction NC10. There are simultaneous choices
$b(v),\rho_v,\psi_v$ satisfying (RC.1), such that
\[
 b(\bar v)=!b(v)\quad(\bar v\ne v),\qquad
 p(s(\bar v),b(\bar v),Q(\bar v),k(\bar v),\psi_{\bar v})
 =p(s(v),b(v),Q(v),k(v),\psi_v).                          \tag{RC.9}
\]
Choose one representative of each orbit of the involution. At that
representative choose a reading and a descent. On a two-element orbit,
use the conjugate reading and the same normalized packet; unnormalize
with the opposite side. NC10 proves the character-pulling identity,
including the inversion needed for the change of side.
At a fixed place keep the original choice; there is no requirement
that a Boolean side equal its own negation. The common families ensure
that the packet remains in precisely the same $i,Q,k$ fiber.
Unpacking it into a descent uses $\operatorname{psiAt}_b(\operatorname{psiAt}_b(\psi))=\psi$.

If $\operatorname{GalOf}(\chi(v),\psi_0(v))$ and
$\iota_v\rho_v=\iota_L$, set
$e_v=\iota_v\circ\operatorname{Frac}(O_{s(v),k(v)}\to O_{\rho_v})$.
Equation (RC.8) proves \texttt{SideRead} for $e_v$.
Applying $\iota_v$ to (RC.1) proves
\texttt{ReadsHecke}$(e_v,\psi_v\circ\operatorname{art},\chi(v))$
using the same Artin map. These two conclusions concern the same
embedding $e_v$, as required in the later rank comparison.

\paragraph{RC10: Frobenius comparison supplies the nonempty-local character tie.}
This criterion does not assume that $D.nd.\mathrm{Loc}$ is empty.
Let $\psi_0(v)$ have \texttt{GalOf}$(\chi(v),\psi_0(v))$ and
\texttt{IsPsi}$(\ell_v,Q(v),\chi(v))$.
For every local index $a$, suppose the already specified character
$\psi_{L,a}=D.ntw.\mathrm{psiL}(a)$ satisfies
\[
 \texttt{IsGalOf}\bigl(D,a,\mathrm{univ},
               \psi_{L,a},\chi(w(a)),0\bigr).            \tag{RC.10}
\]
Then for every $g\in G_E$,
\[
 \operatorname{embK}_a(\psi_0(w(a))(g))
       =\psi_{L,a}(g)\quad\hbox{in }\operatorname{Frac}(O_a),
 \qquad D.nd.Q(a)=Q(w(a)).                              \tag{RC.11}
\]

Let $j_a=\mathrm{iotaQC}_a\circ\mathrm{iotaQ}_a$ on
$\operatorname{Frac}(O_a)$. Its specified \texttt{iota\_LL} law is
$j_a\operatorname{embK}_a=\iota_L$.
The two complex characters
$j_a(\psi_{L,a})$ and $\iota_L(\psi_0(w(a)))$ have open kernels:
this is \texttt{NTw.psiL\_open} for the first and \texttt{GalOf} for
the second. Consequently they are locally constant and continuous,
even though $j_a$ need not be continuous on the fraction field.
Outside the finite set consisting of primes above $\ell_a$ and the
ramification of $\chi(w(a))$, equation (RC.10) with exponent $0$
and the Artin convention give the same value
$\chi(w(a))(P)^{-1}$ at every Frobenius. Apply RC4 to their quotient.
Injectivity of $j_a$ and \texttt{iota\_LL} give the first equality
in (RC.11).

For any prime $P$ and valuation above it, an inertia element is killed
by one side of that equality if and only if it is killed by the other.
The literal \texttt{NTw.psiL\_ram} and RC5 identify the sets where an
inertia element with nontrivial value exists. Their exact
biconditionals therefore give the conductor-set equality in (RC.11).
Now let the normalized packet at $w(a)$ have index $s(w(a))$, side
$b(w(a))$, and support $Q(w(a))$. If
\[
        \mathrm{twOf}(a)=s(w(a)),\qquad
        \mathrm{twSide}(a)=b(w(a)),                       \tag{RC.12}
\]
then (RC.11)--(RC.12) prove every clause of \texttt{LocsCarry}.
This proves the character and conductor portions from the weaker
Frobenius-value interface. It does not construct (RC.10), or choose
the local index and side compatibly with an independently built stage.

\paragraph{RC11: the exact simultaneous choice exists under its stated supplier hypotheses.}
Retain one $D$, a character frame $cf$, $M:\texttt{MtFam}(D)$,
$X:\texttt{NExpTied}(D)$, $F:\texttt{TwFrameC}(D)$, and the global
reciprocity input of RC3. Assume $L_0/\mathbb Q$ is Galois and
$D.nd.\mathrm{Loc}$ is empty. Supply functions
$\xi:\mathrm{Tup}\to\mathrm{HChar}(E)$,
$\chi(v)$, $Q(v)$, $k(v)$ and $s(v)$ with the following exact properties.

\begin{enumerate}
\item Each $s(v)$ has a reading at $v$ on some side, and
$s(\bar v)=s(v)$.
\item Every \texttt{NChoice} $C$ on this same $D$ with $C.Q=Q$
has a \texttt{TupTied} $tp$ with $tp.\xi=\xi$ and $tp.\psi=\chi$.
\item $\chi(v)$ is \texttt{IsPsi} for $\ell_v,Q(v)$, has order dividing
$\ell_v^{k(v)}$, and $L$ contains a primitive root of that order.
\item $Q(v)=\varnothing$ off the nonreference support and
$Q(v)\cap S_0(s(v))=\varnothing$.
\item Equality of the nonnegative-valuation predicates on $L_0$
at $v,v'$ implies $k(v)=k(v')$ and $\chi(v)=\chi(v')$.
Also $\chi(\bar v)=\chi(v)$ and $k(\bar v)=k(v)$.
\item For every $v,b,\psi$, every proof that $(s(v),k(v),Q(v),\psi)$
is admissible, and every $j\in\operatorname{Fin}(D.R.m)$, assume the
following implication. If $v$ is nonreference, $\ell_{s(v)}=\ell_v$,
and \texttt{TwOK} holds, then \emph{for every} $e$ satisfying both
\texttt{SideRead} and \texttt{ReadsHecke} for $\chi(v)$,
the exact predicates \texttt{LevEq} and \texttt{CellWt} of RC2 hold
for $M(\mu_{\mathrm{Tw}}(v),j+1)$ and this cell.
\end{enumerate}

Then there are $C:\texttt{NChoice}(D)$ and
$tp:\texttt{TupTied}(D,cf,C)$ with
\[
 tp.\xi=\xi,\quad tp.\psi=\chi,\quad C.Q=Q,\quad C.k'=k,
 \quad C.\mathrm{twI}=s,                                \tag{RC.13}
\]
all choices have \texttt{TwOK} for their unnormalized character
$C.\psi'\circ\operatorname{art}$, and all six fields of
\texttt{ChoiceSide} hold.

To prove this, a reading first gives $\ell_{s(v)}=\ell_v$.
The exact ramification condition makes $Q(v)$ a function of $\chi(v)$:
if $\chi(v)=\chi(v')$, their support sets are equal. Choose the character
$\psi_0(v)$ of RC3 as a function of
$(\chi(v),\ell_{s(v)}^{k(v)})$ alone. Thus all four families needed
in RC9 are conjugation-invariant. RC5--RC6 give a descent for every
reading, and RC9 chooses coherent $\rho_v,b(v),\psi_v$.
RC8 gives $\iota_v$ and RC9 gives both reading predicates for the
same $e_v$. Hypothesis 6 now applies, and RC2 supplies the exact rank
equalities. NC9 constructs the actual finite DVR extensions,
integral bases and normalized \texttt{NChoice}; it does not assume the
desired integral form isomorphism. Hypothesis 2 gives $tp$ on this $C$.

Here are the six fields of \texttt{ChoiceSide}, including the ones that
remain nonvacuous when the local index set is empty.
The \texttt{NChTie} local index, side and support equalities are vacuous;
its global common-character, conjugation and denominator identities
are (RC.1), (RC.9) and the reading's denominator law.
The six subfields of \texttt{ChoiceC} are index conjugacy,
\texttt{SidePl}, the fraction-field side equality, opposite sides on
two-element orbits, prime-to-$\ell$ power comparison at a common index,
and equality of conjugate packets. All but the power comparison follow
directly from RC8--RC9.

For that comparison, suppose $s(v')=s(v)$. If the sides agree, RC8
identifies the two nonnegative-valuation predicates on $L_0$.
If they differ, conjugate the reading at $v'$; the side then agrees,
and hypothesis 5 together with its conjugation clause again identifies
$k$ and $\chi$. The chosen $\psi_0$ and $Q$ therefore agree.
RC7 gives exactly the required power relation on normalized packets.
The \texttt{off} field is hypothesis 4. The \texttt{adm} field is RC6.

For \texttt{PsiChoiceCompatible}, if
$\iota_L(y)=\chi(v)(P)^{-1}$ at an allowed Frobenius $\sigma_P$,
the Artin convention gives $y=\psi_0(v)(\sigma_P)$.
Equation (RC.1) is then precisely the required comparison in
$\operatorname{Frac}(C.Oe(v))$.
Finally \texttt{eMt\_level} uses the very $\iota_v$ chosen above and
the subgroup/character equalities supplied by hypothesis 6.
Its literal conclusion is equality of levels and character values
under that complex extension. It does not assert that an arbitrary
integral basis isomorphism constructed by NC9 intertwines every Hecke
operator. Thus no extra equivariance is inferred from rank equality.

\paragraph{RC12: prescribed tuple assignments give an actual MuW row datum.}
Fix $B,f,T$ and one supplied $w:\texttt{ColW}(B,f,T)$.
Let $W_t$ be a conjugation-stable set of coefficient places.
Assume an existing $\texttt{RowsRest}(B,w.C)$; SE8 supplies one under
its stated column hypotheses. The exact range law of $w.I.\xi$ contains
every character $T.\xi_{\mathrm{Tw}}(f,v)$, so choose
\[
 w.I.\xi(m_t(v))=T.\xi_{\mathrm{Tw}}(f,v),\qquad
 m_a(v)=
 \begin{cases}m_t(v)&v\in W_t,\\w.C.\mu_0&v\notin W_t.\end{cases}
                                                               \tag{RC.14}
\]
The right side of the first equality is unchanged by $v\mapsto\bar v$:
the residue characteristic and $T.\psi$ are unchanged. Injectivity
of $w.I.\xi$ therefore gives $m_t(\bar v)=m_t(v)$.
Conjugation-stability of $W_t$ gives the same identity for $m_a$.
Moreover $m_a(v)$ is either $\mu_0$ or $m_t(v)$.
Replace precisely the five tuple-assignment fields
$(\mathrm{muAt},\mathrm{muAt}_{\mathrm{conj}},\mathrm{muTw},\mathrm{muTw}_{\mathrm{conj}},\mathrm{muAt}_{\mathrm{or}})$
of the existing row datum by these values. These are the only row-record
laws involving those assignments, so all other fields are retained.
The three \texttt{MuW} equations are exactly (RC.14) and its two cases.

The source's fuller \texttt{reset} also replaces the two weight arrays
by their prescribed integral values and $\mathrm{shaLen}$ by
$\mathrm{shaTrue}$. Those additional interfaces are separate from this
tuple proof. In particular, changing tuple assignments in an existing
row record does not by itself prove that an already supplied geometric
\texttt{RowPk}, its numerical data, or a local arithmetic stage remains
tied to that changed record.

\paragraph{RC13: the actual row lift preserves MuW and supplies all ten MuIsW fields.}
Now retain the actual $w:\texttt{ColW}$,
$P:\texttt{Pre2}(K,w)$ and $x:\texttt{St}(P)$.
The source row lift $q=P.q$ has $W_0$ with \texttt{MuW} and defines
$W=q.p.\mathrm{what}_2(q.e,W_0,q.n,\eta,h_\eta)$.
Its defining equations are
$W.\mathrm{muAt}=W_0.\mathrm{muAt}$ and $W.\mathrm{muTw}=W_0.\mathrm{muTw}$.
The stage $w.\mathrm{st}(x.o.Z)$ uses this exact $W$.
Consequently it satisfies \texttt{MuIsW} with the original $w.I.\xi$,
$T$ and $P.W_t$.

For completeness, the ten fields have the following individual sources.
Injectivity is $w.ok.\xi_{\mathrm{inj}}$; the range is $w.\xi_{\mathrm{range}}$;
the reference value is $w.l.\xi_0$.
The two distinguished values are $w.ok.\xi_{\mathrm{star}}$ and
$w.ok.\xi_{\mathrm{star'}}$, transported through the literal identity
$w.C_{\mathrm{eq}}$ of the reference frame.
The twist, on-support and off-support clauses are $q.\mu$,
preserved by the two defining equations of $W$ above.
The support-residue and off-support character clauses are $P.\mathrm{sub}$
and $P.\mathrm{supp}$. This verifies all ten fields on the actual stage;
no auxiliary zero-form model is used.

The actual frame equality
$cfOf(x)=\operatorname{Fr.at}(w.\mathrm{st}(x.o.Z),f)$
also follows from $w.C_{\mathrm{eq}}$.
The conductor function $QOf(P)$ equals $T.Q$ at every place:
on $P.W_t$ this is its definition; off it, $T.\psi=1$, so the exact
\texttt{IsPsi.ram} law forces $T.Q=\varnothing$.
Thus any $C$ on this same stage with $C.Q=QOf(P)$ satisfies the
support condition of NC16, which supplies a tuple tie with exactly
$tp.\xi=w.I.\xi$ and $tp.\psi=T.\psi$.
The literal \texttt{ColW}, \texttt{Pre2}, \texttt{RowPk} and
\texttt{St} records remain simultaneous inputs here; this transport is
not an existence proof for all their geometric and arithmetic fields.

\paragraph{RC14: a rank-one conjugation detects the parahoric scalar-action subgroup.}
At a split real-field prime, put $K=F_v$ and let $\mathcal O$ be its
valuation ring. The unitary group there is $\operatorname{GL}_n(K)$:
under $E\otimes_F K=K\times K$, a unitary matrix is the pair
$(h,h^{-T})$. Thus every $h$ has a unique adelic lift supported at $v$.
It is integral exactly when $h\in\operatorname{GL}_n(\mathcal O)$.
This lift is an actual finite adele: its two components at $v$ are
prescribed, it is the identity elsewhere, and expansion in a fixed
$F$-basis of $E$ gives finite-adele coefficients and inverses.

Let $P_r\le\operatorname{GL}_n(\mathcal O)$ be the inverse image of
the upper block parabolic of block sizes $r,n-r$.
Suppose a nonzero function $F$ on $G_f$ has the prescribed scalar
action on $P_r$, and its scalar character has the same kernel as
$u\mapsto\xi(\det u_{22})$. Suppose $\xi(a)\ne1$ for some
$a\in\mathcal O^\times$. Every integral element supported at $v$ that
acts by a scalar on $F$ belongs to $P_r$.

Indeed, write its local matrix as $h$. If $h\notin P_r$, there are
$i\ge r$, $j<r$ with $h_{ij}\in\mathcal O^\times$.
Set $t=(a-1)/h_{ij}\in\mathcal O$, and define
\[
 A=I+t(he_j)e_i^T,\qquad
 B=I+t e_j(e_i^Th)=h^{-1}Ah.                            \tag{RC.15}
\]
The identity $\det(I+uv^T)=1+v^Tu$ gives
$\det A=\det B=a$. Both matrices and their inverses are integral:
their entries are integral, their determinant is a unit, and the
adjugate formula gives integral inverses.
Matrix $A$ changes only column $i\ge r$; matrix $B$ changes only row
$j<r$. Both therefore lie in $P_r$. Their lower-right block
determinants are respectively $a$ and $1$.
Their adelic lifts $u,w$ satisfy $w=\kappa^{-1}u\kappa$.
The scalar for $u$ on $F$ is nontrivial and that for $w$ is $1$.
The scalar for $\kappa$ is nonzero, since right translation is invertible
and $F\ne0$. Applying the two equal products $\kappa w=u\kappa$
to the nonzero line spanned by $F$ forces those two scalars to agree,
a contradiction. Only equality of character kernels was used, so this
also applies to two complex readings of one coefficient character.

\paragraph{RC15: an automorphism cannot send one parahoric into both opposite parahorics.}
Let $0<r<n$. At a fixed split prime, the two possible embeddings
$E\hookrightarrow F_v$ are conjugate. One local matrix is consequently
the inverse transpose of the other. In the compact group
$K_v=\operatorname{GL}_n(\mathcal O)$, their parahorics are the inverse
images $P_+,P_-$ of opposite block parabolics in
$\operatorname{GL}_n(\kappa_v)$.
They have equal finite index, since inverse transpose interchanges
them. Their intersection has strictly larger index than either one:
$I+E_{0r}$ belongs to $P_+$ and not $P_-$.

Suppose $\theta$ is the source's \texttt{FrameAut}: it preserves the
global integral group and preserves support at each prime, in both
directions. Its restriction to the subgroup supported at $v$ is
therefore an automorphism of $K_v$. If
\[
                  \theta(P_\pm)\subseteq P_+\cap P_-,
                                                               \tag{RC.16}
\]
then preservation of finite subgroup index gives
$[K_v:P_\pm]\ge[K_v:P_+\cap P_-]>[K_v:P_\pm]$, impossible.
Consequently, if the image under $\theta$ of a parahoric for some
split embedding lies in the parahorics for two split embeddings at
the same prime, those two embeddings are equal.
There are only two possibilities because $E/F$ is quadratic:
write $E=F\oplus Fy$ with $\bar y=-y$ and $y\ne0$; the image of $y$
is determined by $y^2\in F$ up to sign. This proves the literal
split-place equality used in \texttt{LevCmp.levEq}, including its
orientation.

\paragraph{RC16: the common scalar action recovers the additive level character.}
Let the residue characteristic at the split prime be $p$.
Let $U_r\le K_v$ be the inverse image of the subgroup
\[
 \left\{\begin{pmatrix}u&B\\0&D\end{pmatrix}:
 u\text{ upper unitriangular of size }r,\
 D\in\operatorname{GL}_{n-r}(\kappa_v)\right\},
 \qquad 2\le r\le n.
\]
Its index in $K_v$ is finite and prime to $p$.
For a direct proof, it contains the inverse image of the full upper
unitriangular subgroup $U_n(\kappa_v)$. Reduction
$\operatorname{GL}_n(\mathcal O)\to\operatorname{GL}_n(\kappa_v)$
is surjective by entrywise lifting and the unit-determinant test.
The index of that smaller subgroup is
\[
 \frac{|\operatorname{GL}_n(\kappa_v)|}{|U_n(\kappa_v)|}
       =\prod_{a=1}^{n}(q_v^a-1),                       \tag{RC.17}
\]
since an ordered basis has
$\prod_{a=0}^{n-1}(q_v^n-q_v^a)$ choices and $U_n$ has
$q_v^{n(n-1)/2}$ elements. Every factor on the right is prime to $p$,
and the index of $U_r$ divides it.

Let nonzero functions $F,F'$ transform on $U_r$ by characters whose
additive factors are $\alpha,\alpha'$ and whose lower-block determinant
factors are arbitrary. Assume $\alpha,\alpha'$ are trivial on the
maximal ideal. Let $\theta$ preserve $K_v$ and suppose that, for every
$u\in U_r$, $\theta(u)$ acts by scalars $c,c'$ on $F,F'$, with
$c^m=(c')^m$, where $\gcd(m,p)=1$.
Then $\alpha(y)=\alpha'(y)$ for every $y\in\mathcal O$.

To prove this, let $H\le K_v$ consist of elements acting by a scalar
on both functions. A nonzero function has a unique nonzero scalar
for any such element; these scalars define homomorphisms
$\lambda,\lambda':H\to\mathbb C^\times$.
Set $\eta=(\lambda/\lambda')^m$.
The hypotheses give $\theta(U_r)\le\ker\eta\le H\le K_v$.
Thus $d=[H:\ker\eta]$ is finite and divides
$[K_v:\theta(U_r)]=[K_v:U_r]$, so $\gcd(d,p)=1$.
For the transvection $x_y=I+yE_{01}$, $y\in\mathcal O$,
the lower-block determinant is $1$ and the top superdiagonal sum is $y$.
It lies in both original level groups and in $H$, and hence
\[
       (\alpha(y)/\alpha'(y))^{md}=1.
                                                               \tag{RC.18}
\]
Also $p y$ belongs to the maximal ideal, so
$\alpha(y)^p=\alpha'(y)^p=1$. Bezout for $p$ and $md$ now gives
$\alpha(y)/\alpha'(y)=1$. The source uses $m=p-1$.
The conclusion concerns values on integral $y$, which are exactly the
superdiagonal values occurring in the level-character formula.

\paragraph{RC17: one explicit integral frame gives a trivial stabilizer class.}
Let $n\ge2$, $r\ge2$, and let a split prime $P\mid p$ satisfy
$p\ge2n$ and $p\notin P^2$. Then an adelic element $g$ can be chosen
such that
\[
 \gamma\in U_n(E),\quad
 g^{-1}\gamma g\in U_{\max},\quad
 (g^{-1}\gamma g)_P\in U_r
             \quad\Longrightarrow\quad\gamma=1.         \tag{RC.19}
\]
Use the integral matrix $A$ with first column $(1,\ldots,1)^T$,
second column $(0,1,\ldots,n-1)^T$, and remaining columns the corresponding
standard basis columns. It is lower triangular with diagonal entries
$1$, so $A\in\operatorname{GL}_n(\mathbb Z)$. Lift it at the chosen
split prime and use the identity at all other primes. This gives
$g\in U_{\max}$ by the explicit split lift of RC14.

Since $g$ is integral, the hypotheses make $\gamma$ an everywhere
integral rational unitary matrix. Each of its entries is an algebraic
integer of absolute value at most $1$ under every complex embedding:
the sum of squared absolute values in its row is $1$.
Such a nonzero entry is a root of unity. One elementary proof is to
consider its powers: their conjugates remain bounded, so the integral
coefficients of their monic characteristic polynomials lie in a fixed
finite set. Only finitely many roots of those finitely many polynomials
occur, forcing two powers to agree.
At least one row entry is nonzero; its absolute value is $1$, so every
other entry in that row is zero. Thus row $i$ has one nonzero entry
$\zeta_i$, in column $\sigma(i)$, where $\zeta_i$ is a root of unity.

Reduction of $\gamma A=AH$, with $H\in U_r$, on the first column gives
$\zeta_i\equiv1\pmod P$.
A torsion unit congruent to $1$ at this unramified odd prime is $1$.
Here is the exact elementary argument. If a nontrivial such unit has
finite order, a power $u=1+z$, $z\in P$, has prime order $a$.
The geometric sum $1+u+\cdots+u^{a-1}=0$ shows $a\in P$,
so $a=p$. Expanding it modulo $P^2$ gives
\[
       0=p+\frac{p(p-1)}2 z+z^2 t
\]
for some integral $t$. Since $p$ is odd and $p,z\in P$, this forces
$p\in P^2$, contrary to the hypothesis. Therefore all $\zeta_i=1$.

The second column now gives
$\sigma(i)-i\equiv H_{01}\pmod p$ for every $i$.
The difference between any two integers $\sigma(i)-i$ has absolute
value at most $2(n-1)<p$, so divisibility by $p$ forces exact equality.
Their common value is nonnegative from $i=0$ and nonpositive from
$i=n-1$, hence is zero. Thus $\sigma(i)=i$ and $\gamma=1$.

For the actual $D$, the source hypotheses supply these inequalities
and unramifiedness. Put $n=D.R.nV=r+t$.
The congruence $g_k\equiv1\pmod{Q_Sqd}$ and primality of $g_k$ imply
$g_k\ge Q_Sqd\ge q>4tr+1\ge2(r+t)$.
The specified count of primes above $g_k$ is $[E:\mathbb Q]$.
In the equality $\sum_{P\mid g_k}e_Pf_P=[E:\mathbb Q]$, every summand
is at least $1$, so all $e_P=f_P=1$. In particular $g_k\notin P^2$.
This uses the actual type-prime fields; it does not assert existence
of those primes without the host/type input.

\paragraph{RC18: the stabilizer class gives a nonzero form for every nonzero weight.}
Let $J$ have the source's product-level shape: integral parahoric
conditions at the ramified split primes, integral unipotent conditions
at the nonempty set of type primes, and the integral maximal compact
condition outside those sets. A product level permits the component
of any $\kappa\in J$ at one prime to be isolated inside $J$.
Each isolated component is integral by its stated local condition;
the complement condition treats all other primes. Consequently
$J\subseteq U_{\max}$, and every element of $J$ satisfies the specified
type-prime unipotent condition.

Choose the class $g$ of RC17 at one actual type prime. Its rational
stabilizer $\Gamma\cap gJg^{-1}$ is trivial.
For a nonzero $w$ in any coefficient representation $W$ and any
$J$-character $\chi$, define
\[
 F(\gamma g\kappa)=\chi(\kappa)\rho(\gamma)w
       \quad(\gamma\in\Gamma,\ \kappa\in J),
 \qquad F=0\text{ off }\Gamma gJ.                        \tag{RC.20}
\]
The expression is unique: if
$\gamma g\kappa=\gamma' g\kappa'$, then
$\gamma'^{-1}\gamma\in\Gamma\cap gJg^{-1}$, hence both factors agree.
Multiplication on the left by $\Gamma$ and on the right by $J$
preserves the double coset, so (RC.20) satisfies both transformation
laws on it and on its complement. Also $F(g)=w\ne0$.
For an irreducible weight the space is nonzero by the definition of
irreducibility. This proves the nonzero-form conclusion used for both
the complex \texttt{MtReal} model and the coefficient-field cell.
The congruence-kernel field in \texttt{MtReal} converts the former
function into its smooth form-space range, and the cell's actual
\texttt{real} isomorphism supplies its preimage in \texttt{MtV}.

\paragraph{RC19: the local arguments give the literal global LevEq.}
Retain the exact hypotheses of \texttt{LevCmp.levEq} on
$m:\texttt{MtReal}$ and the cell $c$ of $X:\texttt{NExpTied}$:
the actual type-prime data \texttt{GPlaces};
product, ramified-parahoric, type-unipotent and outside-compact
descriptions for $m$; the ramified and type descriptions for $c$;
$\mathrm{badV}\subseteq gPl\cup qPl$;
nontriviality of the common ramified component on some local unit;
and the exact \texttt{Joint} datum.
The cell's type-character description must hold for every complex
reading extending the specified side. Put $S=gPl\cup qPl$.
Then
\[
             \texttt{LevEq}(m.J,m.\chi,c.J,e\,c.\chi).
                                                               \tag{RC.21}
\]

We spell out the role of \texttt{Joint}. It gives an automorphism
$\theta$ preserving integrality and each prime support, a subgroup $J_1$,
another side embedding $e'$, and an injective $e'$-semilinear map
$\omega$ from the cell's weight space to a complex vector space.
The subgroup contains specified parahorics over $qPl$ and specified
unipotent groups at all type primes. Every element of $J_1$ acts,
after applying $\theta$, by scalars on all the realized $m$-forms and
on all cell forms after $\omega$; at a type prime their
$(g_k-1)$st powers agree. This is substantially more than equality
of ranks or an arbitrary linear isomorphism of two form spaces.

The cell is a product level directly from its \texttt{J\_mem} formula.
Indeed, isolate one adelic component using the idempotent of that
prime; every listed matrix condition becomes either the original
condition or the identity condition. Integrality is preserved.
Its outside-compact description follows from its \texttt{hyp} field,
the bad-set inclusion, and injectivity of $e$, which preserves the
character kernel. RC18 supplies nonzero forms on both sides.
Their preimages under the actual realization maps, and injectivity
of $\omega$, give two nonzero complex-valued weight functions to
which RC14--RC16 apply.

At a type prime, two local multiplicative components representing the
same Hecke character are equal. To see uniqueness without a continuity
shortcut, both components kill sufficiently deep principal units.
Approximate an arbitrary local unit by a nonzero element of $F$ so that
the quotient lies in both kernels. Their prescribed values on that
rational element agree, and therefore so do their values on the local
unit. This works for every element of $F_v^\times$, with the depth of
approximation adjusted to its valuation. RC16, with $p=g_k$ and
$m=g_k-1$, identifies the additive factors on integral elements.
The groups have the same unipotent matrix condition. The lower block
of an integral parahoric matrix has unit determinant: reduction is
block upper triangular and the full determinant is nonzero.
Its superdiagonal sum is integral. The character formulas consequently
agree on every supported group element.

At a ramified prime, \texttt{Joint} provides a parahoric at some split
embedding. RC14, applied to both nonzero functions, puts its image
under $\theta$ inside both proposed parahorics. The change from $e$
to $e'$ has no effect on the coefficient character's kernel.
RC15 forces the two split embeddings to be identical, and uniqueness
of the common Hecke component gives equality of their determinant
characters. Thus the two levels and characters agree on each
single-prime subgroup over $S$.

Finally, $S$ is finite: the type set lies over two rational primes,
and a Hecke character is ramified at only finitely many finite primes.
For the latter assertion, continuity and a complex neighborhood
containing no nontrivial subgroup force the character to kill the
unit factors outside a finite set.
For $\kappa\in m.J$, isolate its components at the finitely many primes
of $S$, successively divide them off, and leave a factor supported
outside $S$. The isolated factors belong to $c.J$ with equal character
values by the local comparison; the remaining factor belongs by the
common outside-compact rule and has character $1$ on both sides.
Multiplying proves membership and character equality globally.
Reverse the roles of the two product levels to obtain equality of
subgroups. This proves (RC.21), and RC2 can then be applied with the
separate exact \texttt{CellWt} input.

\paragraph{RC20: an actual coherent index and side selector on the twisted places.}
Assume $L_0/\mathbb Q$ is Galois, the literal
\texttt{TwCoverPair} holds, and an index $i_0\in D.\mathrm{ntw}.Tw$
is supplied. Write $v\sim v'$ if the two valuations have the same
positive-valuation predicate on $\mathcal O_{L_0}$. Write $v\approx v'$
if $v\sim v'$ or $\bar v\sim v'$.
Conjugation preserves $\sim$: normality of $L_0$ lets complex conjugation
act on $L_0$, and after clearing an integer denominator it acts on
its ring of integers as well. Thus $\approx$ is an equivalence relation,
and $v\approx\bar v$.
The coverage hypothesis says precisely that the prime of
$\mathcal O_{L_0}$ below every \emph{twisted} coefficient place is
the center of one of the two side embeddings of some actual twist index.
Consequently every twisted $v$ admits some \texttt{ReadsAt}$(i,b,v)$.

For each $\approx$-class choose an index satisfying
\[
 \exists u,b,\quad v\approx u\quad\hbox{and}\quad
                  \texttt{ReadsAt}(i,b,u),                    \tag{RC.22}
\]
if one exists, using $i_0$ otherwise. Equivalence makes the set of
eligible indices exactly the same throughout a class, so a fixed
choice function gives $s(v')=s(v)$ whenever $v\approx v'$.
At a twisted $v$, (RC.22) is nonempty. A reading at $u$ transports
along $\sim$ without changing its side; if conjugation is needed,
conjugating the reading reverses the side. Hence choose a preliminary
$b(v)$ with $\texttt{ReadsAt}(s(v),b(v),v)$ at every twisted $v$.
Outside that set its value may be arbitrary.

Choose one representative $o(v)\in\{v,\bar v\}$ of each conjugation
orbit. Define the final side by $b(v)$ at the representative and
by the opposite of the representative's side at the other place.
The reading conjugation identity and $s(\bar v)=s(v)$ show that the
final side still reads the selected index. At a two-element orbit
the sides are opposite; at a fixed place no impossible sign reversal
is required. This is exactly \texttt{exists\_settings}.
Its reading conclusion holds on the twisted places only.
It does not give the all-place hypothesis \texttt{hs} of RC11,
nor does it supply a reading for an empty index type.

\paragraph{RC21: finite integral bases and prime avoidance give the actual ExistsData.}
Fix the exact expansion datum $\mathcal E$, the function $N_0$ on
twist indices, and the following five source hypotheses:
$\texttt{aux\_3}(D,\mathrm{bk0},\mathcal E,N_0)$,
\texttt{TwistedPairSettingHyps}, \texttt{FullRank},
\texttt{adm7}, and \texttt{TwCoverPair}, together with
$L_0/\mathbb Q$ Galois. The first is the actual twisting theorem
at every nonzero index: given its setting, any positive error budget,
the full-rank theta datum, any finite integral theta basis, and any
finite forbidden set, it supplies a packet satisfying
\texttt{GoodTw}. We use that hypothesis with its full meaning,
including the lattice-index inequalities, the finite-extension
quantifiers in \texttt{IdxAt}, and the three arithmetic parts
\texttt{partAT}, \texttt{partC}, \texttt{partD}.
We do not prove these inequalities by choosing an arbitrary character.
Under these hypotheses the literal conclusion is
\[
 \exists S\ \ \forall A:Tw\to\operatorname{Finset}(Pl(E))\
 \exists(t,sd)\quad \texttt{GoodData}(S,t,sd),             \tag{RC.23}
\]
with $Q(t(v))\cap A(s(v))=\varnothing$, with $t$ constant on
$\approx$-classes of twisted places, and, if the twisted-place
set is finite, with the rational-prime supports of packets from
different $\approx$-classes pairwise disjoint.

Here is the integral basis construction, which occurs \emph{before}
the forbidden sets $A$ are specified. Fix an active $i$ and
$1\le j\le D.R.m$, and write $O=O_i$.
At the trivial twist $k=0,Q=\varnothing,\psi=1$, all seven
\texttt{TwOK} clauses hold. The hypothesis \texttt{adm7}
(literally \texttt{admT4}) places the scalar extensions
$1\otimes\theta(a)$, for $a\in\mathcal E.\mathrm{Dat0}(i,j,b)$,
inside the image of the given finitely generated module
$\mathrm{MtLat}$ under $\mathrm{thP}$.
The finite injective local extension $O\to O_{i,0}$ is flat, since it
is torsion-free over the DVR $O$, and faithfully flat, since it is local.
Thus $x\mapsto1\otimes x$ is injective. The containing image is a
finitely generated $O$-module by finiteness of $O_{i,0}/O$.
Noetherianity of $O$ and injectivity show that the original theta span
is finitely generated.

Its image under the actual integral averaging operator
$AO=|G_0|e_O$ is also finitely generated and torsion-free.
It has a basis chosen from the vectors $AO\,\theta(a)$ themselves.
Indeed, first express finitely many module generators as finite
combinations of this generating family; their combined supports give
a finite generating subset. Choose one with minimal cardinality.
If it has a nontrivial relation, choose a nonzero coefficient of
least DVR valuation. It divides every other coefficient. Divide the
relation by it using torsion-freeness. The chosen coefficient is now
$1$, so its vector is a combination of the others, contrary to
minimality. Hence the chosen subset is a basis. Together with the
retained \texttt{FullRank} identity this gives all fields of
\texttt{TwIn.Ok}. Use empty subsets at inactive indices or out-of-range
$j$. Choice over $(i,j,b)$ gives one family $S$ independent of $A$.

We next prove the required finite avoidance step. For a finite set
$B$ of primes of $E$, let $B^+$ be the union, over $P'\in B$,
of all primes containing the nonzero integer $\mathrm N(P')$.
This set is finite, by factorization of nonzero ideals in
$\mathcal O_E$. If $P'$ lies over the rational prime $p$, its norm is
$p^f$ for some $f>0$; every prime $P$ over $p$ therefore contains
$\mathrm N(P')$. It follows that $Q\cap B^+=\varnothing$ implies
that no prime in $Q$ lies over the same rational prime as a prime in
$B$. This is the exact symmetric relation \texttt{Apart}.

For each active $i$ and each finite $B$, apply \texttt{aux\_3}
with the fixed basis $S_i$, the actual positive budget
$\mathrm{twEps}(D,i)$, its supplied setting, and forbidden set
\[
             \mathrm{AvRec}(D)\ \cup\ A(i)\ \cup\ B^+ .
                                                               \tag{RC.24}
\]
The resulting $(Q,k,\psi)$ has the same \texttt{GoodTw} properties
after forgetting the extra forbidden sets, is disjoint from $A(i)$,
and is \texttt{Apart} from $B$.
At an inactive index choose $(\varnothing,0,1)$: the required
goodness assertion is conditional on activity and hence vacuous.

An elementary induction now selects good packets at all indices
which are pairwise \texttt{Apart} on any prescribed finite subset
$T\subseteq Tw$. Start with a good packet at every index.
When adding $a$ to a previously treated finite set $T$, keep all
other packets and choose the packet at $a$ with $B$ equal to the
union of their conductor sets. The preceding construction and
symmetry of \texttt{Apart} give every new pairwise condition.
All old conditions are unchanged.

If there are no twisted places, (RC.23) is immediate, with empty
packet function and arbitrary side. Otherwise a single covered
twisted place supplies the index $i_0$ needed in RC20.
Use its selector $s,sd$. If the twisted-place set is finite, take
$T$ to be its finite image under $s$; otherwise take $T=\varnothing$,
as no pairwise-apart conclusion is then required.
Select packets $f(i)$ by the preceding induction and set
$t(v)=(s(v),f(s(v)))$.
A reading at a twisted place makes its index active, so every
required \texttt{GoodPack} clause holds there.
RC20 gives reading, side reversal, conjugation invariance, and
constancy of $t$ on $\approx$-classes.
Finally, if two twisted places have the same selected index,
their readings imply $v\approx v'$: equal sides have the same
center on $\mathcal O_{L_0}$, and opposite sides differ by conjugation.
Consequently distinct $\approx$-classes use distinct indices, and
the finite selection supplies the promised \texttt{Apart} property.
This proves every clause and the quantifier order of
\texttt{ExistsData}.

\medskip
\noindent\textbf{Exact boundary.}
The actual rank, character descent, finite reading choices, local
level comparison and tuple compatibility have now been proved
under the displayed inputs. The simultaneous construction still
requires the supplied global Artin reciprocity map; the actual
ring-class universal property and twist coverage; the exact
geometric and representation-theoretic realization data producing
\texttt{Joint}, local shape, and \texttt{CellWt}; and the genuine
twisting/full-rank/integral-inclusion hypotheses used in RC21.
For nonempty local data, RC10 gives a sufficient exact
Frobenius-character criterion, while the local index and side
identifications remain necessary.
For tuples, the actual common \texttt{ColW}, \texttt{Pre2},
\texttt{RowPk} and \texttt{St} objects are retained; resetting tuple
maps in an unrelated row model does not supply them.
The bare stage with $Tw=Loc=\varnothing$ supplies none of the
all-place coverage required for RC11.
The final theorem \texttt{choices\_exist\_for} still depends on these
simultaneous suppliers and their stated analytic hypotheses.
